% Einheit 3 von 3 – Wurzelgesetze und rationale Exponenten (bank/reelle-zahlen/e3.jsonl, zusammenbau v0.1)
%% TODO Zweigzeile Teil 1: was hier gelernt wird (Satz in Schülersprache aus dem Einheitstitel) – Einheit 3
\einheitenkopf[e3]{Einheit 3 von 3 · Wurzelgesetze und rationale Exponenten}
\zweigzeile{ab Kl. 9 (am Gymnasium ab Kl. 8) · keine P10-Aufgabe}
\setcounter{aufgabe}{16}

%% TODO Titel: Ich-kann-Satz fehlt in der Bank (Platzhalter: Kettenname) – Nr. 17
%% TODO Anweisung: ein Satz über der ersten Teilaufgabe fehlt in der Bank – Nr. 17
\begin{aufgabe}{Wurzelgesetze und rationale Exponenten}
%% TODO \rechenplatz{n}: Schrittzahl des Grundfalls fehlt in der Bank (nur bei fester Schreibform, 2.3 b)
\begin{teile}
\teil $\sqrt{4 \cdot 81}$ – was steht unter der Wurzel? \\ \kreuz{Produkt oder Quotient: trennbar} \\ \kreuz{Summe oder Differenz: erst ausrechnen}
\teil $\sqrt{9 \cdot 64}$ – erst malnehmen, dann Wurzel; und die Wurzeln einzeln. Gleiches Ergebnis? \leerfeld und \leerfeld
\teil $\sqrt{\frac{64}{9}}$ – als Bruch? \leerfeld
\teil $\sqrt{28}$ – teilweise Wurzel ziehen? \leerfeld
\teil $4\sqrt{3} + 2\sqrt{3}$ – zusammengefasst? \leerfeld
\teil $\sqrt{81 + 144}$ – Ergebnis? \leerfeld
\teil $81^{\frac{1}{2}}$ – als Wurzel und als Zahl? \leerfeld
\teil $\sqrt[3]{64}$ – Ergebnis? \leerfeld
\teil $5^{\frac{1}{2}} \cdot 5^{\frac{1}{2}}$ – Ergebnis? \leerfeld
\teil $\frac{1}{\sqrt{3}}$ – mit rationalem Nenner? \leerfeld
\teil Begründe mit den Potenzgesetzen: $\sqrt{a} \cdot \sqrt{b} = \sqrt{a \cdot b}$ für $a, b \ge 0$.
\teil $x^3 = 512$ – Lösung? $x =$ \leerfeld
\teil $\sqrt{18} + \sqrt{8}$ – vereinfacht? Schreibe das Ergebnis auch mit einer Potenz von $2$.
\end{teile}
\end{aufgabe}

%% TODO Titel: Ich-kann-Satz fehlt in der Bank (Platzhalter: Kettenname) – Nr. 18
%% TODO Anweisung: ein Satz über der ersten Teilaufgabe fehlt in der Bank – Nr. 18
\begin{aufgabe}{Wurzelgesetze und rationale Exponenten – Fehler finden, Begründen, Darstellungswechsel, Anwendung}
\begin{teile}
\teil Finde den Fehler und rechne richtig: \rechnung{\sqrt{81} + \sqrt{144} &= \sqrt{225} \\ &= 15}
\teil Zeige mit einem Zahlenbeispiel, dass $\sqrt{a + b}$ im Allgemeinen nicht $\sqrt{a} + \sqrt{b}$ ist.
\teil $\sqrt[5]{x^3}$ – als Potenz? \leerfeld
\teil Ein quadratisches Grundstück hat $800\,\text{m}^2$. Gib die Seitenlänge exakt und möglichst einfach sowie gerundet an.
\end{teile}
\end{aufgabe}
